\section{Related Work}
\label{sec:related}

\para{Credit assignment in cooperative MARL.}
Credit assignment---determining each agent's contribution to a shared team reward---is a central challenge in cooperative \marl. Value decomposition methods factorize the joint action-value function into individual agent utilities. \qmix \citep{rashid2018qmix} uses a monotonic mixing network conditioned on global state, while Weighted \qmix \citep{rashid2020weighted} relaxes the monotonicity constraint for better expressiveness. QPLEX \citep{wang2021qplex} achieves full expressiveness through a duplex dueling architecture. \coma \citep{foerster2018coma} takes a different approach, using counterfactual baselines to assess each agent's marginal contribution. LICA \citep{zhou2020lica} learns implicit credit assignment via hypernetwork-based mixing critics. These methods address the \emph{agent} dimension of credit assignment but treat the temporal dimension uniformly.

\para{Temporal credit assignment.}
Assigning credit across time---determining which past actions led to current rewards---has been studied primarily in single-agent settings. \rudder \citep{arjona2019rudder} redistributes episodic returns to create return-equivalent MDPs with zero expected future rewards, enabling faster learning with delayed rewards. Hindsight credit assignment \citep{harutyunyan2019hindsight} uses future information to improve temporal credit via hindsight-conditioned policy gradients. More recently, \tartwo \citep{kapoor2024tar2} provides the first principled framework for \emph{joint} agent-temporal credit decomposition, using alternating temporal and agent attention modules to decompose episodic returns into per-agent, per-timestep rewards. \stas \citep{chen2024stas} uses a dual transformer architecture with causal temporal attention and Shapley-based spatial attention for simultaneous decomposition. Both methods assume synchronous agents with uniform decision frequencies. Our work investigates whether the temporal dimension of credit assignment can emerge from heterogeneous horizons rather than being imposed architecturally.

\para{Hierarchical and multi-timescale RL.}
Hierarchical RL explicitly introduces multi-timescale decision-making. \feudal Networks \citep{vezhnevets2017feudal} use a Manager operating at a coarser timescale to set goals for a Worker executing primitive actions. The Option-Critic architecture \citep{bacon2017option} learns temporally extended options and their termination conditions end-to-end. MAVEN \citep{mahajan2019maven} conditions on shared latent variables for coordinated exploration, while RODE \citep{wang2021rode} decomposes the joint action space through learned agent roles. These methods \emph{impose} hierarchical structure by design. We ask whether analogous structure \emph{emerges} when agents simply operate at different frequencies.

\para{Asynchronous multi-agent decision-making.}
The MacDec-POMDP framework \citep{amato2019modeling} formalizes asynchronous multi-agent settings where agents execute macro-actions of different durations. MacroMARL provides environments with native support for heterogeneous macro-actions but does not study credit assignment patterns. \mappo \citep{yu2022mappo} demonstrates that properly tuned PPO with a shared critic achieves strong results in cooperative \marl, but operates synchronously. Our work bridges the gap between asynchronous execution and credit assignment analysis, using the simplest possible algorithm (\ippo) to isolate the effect of temporal heterogeneity.
